\chapter*{Conclusion and Future Work}
\addcontentsline{toc}{chapter}{Conclusion and Future Work}

\section*{Conclusion}

Breakly delivers its central idea: entertainment inside a break that has a
clear end. All Must and Should requirements of the current version are
implemented. A break starts in two taps from an illustrated room, the
countdown stays correct in the background, and every service has a limit or a
natural end: music stops with the break, the TV stops after ten videos, and
each daily game can be played once a day. The app works without an account,
and signing in keeps sessions, streaks and preferences across devices; records
made offline are uploaded later without duplicates. During development the
team replaced the first video source, a short curated list, with the real
YouTube website and a video limit, added the Game Corner, and extended the
radio with eight stations, offline MP3 downloads, a seek bar and lock-screen
controls on iOS.

\section*{Limitations}

The acceptance tests on the iPhone are still to be completed, and the app has
no automated UI tests. The deployment workflow does not yet publish the
backend to a server. The TV and the Game Corner depend on the page structure
of YouTube and NYT Games, which may change, and the music catalogue depends on
the availability of the Jamendo API.

\section*{Future Work}

\begin{itemize}[leftmargin=1.8em,itemsep=2pt]
  \item Complete the acceptance tests and the accessibility test of NFR-08, and
        add automated tests for the session store and the API client.
  \item Deploy the backend with Docker Compose and HTTPS on a VPS and
        automate it with the deployment workflow.
  \item Build the onboarding carousel and password reset designed in Figma.
  \item Use an audio playlist so that the lock screen also offers next and
        previous track buttons, and respect the system Reduce Motion setting.
  \item Add more Spaces and games, and validate the personas and the effect
        on break length with intended users.
\end{itemize}
